\section{What is added, and what is deliberately not repeated}

\subsection{A revision-pinned boundary}

This report was prepared against Forge commit
\code{c98e47c5f804e92880fc1d0e378c1b95832b685c}, not the earlier
\code{6745210} design-round snapshot. The current README describes a 111-page
merged article and eighteen contributing proposals, including an extension
round. Its later merge records elaboration of sixteen core-only Lean files;
that progress must not be erased by repeating the original proposals' older
``Lean unavailable'' status. It still does not install a Forge tactic or report
comparative tactic performance~\cite{forge}.

The review used the current README, merged provenance ledger, baseline section,
scalar nonlinear section, the principal closure sections, and the
Leant/Djex review. It also inspected the current Leant and Djex entry documents
at the revisions in Table~\ref{tab:pins}. This was a targeted source audit,
not an exhaustive line-by-line review of every frozen submission. ``Additional
to Forge'' below means additional to the mechanisms identified in that inspected
merged baseline. It is not a claim of new mathematics or a priority claim over
the noncommutative-algebra literature.

\begin{table}[ht]
\small\centering
\begin{tabularx}{\textwidth}{@{}lX@{}}
\toprule Project & Inspected revision \\\midrule
Forge & \code{c98e47c5f804e92880fc1d0e378c1b95832b685c} \\
Leant & \code{6bf05ad78c467989e68290f2d08bbed40802d485} \\
Djex & \code{e8778f4ebd63e1f9b9b410fa4de8d14a8a04c9e5} \\
\bottomrule
\end{tabularx}
\caption{Repository read snapshots. Mathlib documentation was consulted live;
this is not a tested joint dependency lockfile.}\label{tab:pins}
\end{table}

\subsection{The capability delta}

\begin{longtable}{@{}p{0.28\textwidth}p{0.32\textwidth}p{0.33\textwidth}@{}}
\toprule Existing ingredient & Not a new contribution here & Actual addition \\\midrule\endhead
Commutative ideal certificates & Tracking polynomial multipliers and checking
an expanded identity & Two-sided word contexts $u r v$; exact degree-complete
homogeneous consequences and matrix witnesses. \\
Scalar quadratic positivity & Exact rational $LDL^{\mathsf T}$ as a numerical
linear-algebra method & The unique word Gram matrix for homogeneous degree
$2d$, with $m^d$ words and no commutative coefficient collisions. \\
Scalar cone certificates & Nonnegative combinations of preselected squares &
Adjoint sandwiches $q^*gq$, equality-space projection, and a structural ray
generator for involutive operators. \\
Ore transport & Noncommutativity itself & General free-word relations, not a
single prescribed shift/coefficient commutation law. \\
Matrix observables and tensor lifts & Matrices as finite state-transition data &
Symbolic matrix/operator inequalities valid independently of matrix dimension. \\
General certificate architecture & Opaque receipts, process separation, exact
source identity, and status distinctions & Interpretation-specific replay:
algebra equality, operator positivity, trace positivity, and algebra countermodel. \\
\bottomrule
\end{longtable}

The previous scalar acceptance language uses commuting real indeterminates and
products of nonnegative constraints; the existing closure and Ore mechanisms
have their own explicit semantic fragments~\cite{forge-nonlinear,forge-closure}.
Neither interface should be silently widened to the new one. In particular,
positive semidefinite matrices cannot be multiplied as though they were
nonnegative real numbers.

\subsection{A stronger baseline than bare normalization}

Mathlib already has \code{noncomm\_ring}. It expands noncommutative ring
expressions using simplification and additive normalization, accepts supplied
rewrite lemmas and local hypotheses, and documents limitations around fully
expanded powers~\cite{noncomm-ring}. Therefore this proposal is not ``add a
noncommutative normalizer.'' Its search constructs \emph{new combinations of
hypotheses in left and right contexts}, and its other outputs include
countermodels and positivity receipts. Some worked equalities may already be
easy for \code{noncomm\_ring [*]}, \code{simp}, or \code{grind}. They are
replay and integration examples, not evidence that existing tactics fail.

The current \code{grind} remains a cooperating collection of proof-producing
engines, not a trivial baseline~\cite{grind}. A fair later evaluation needs
\code{grind}, \code{noncomm\_ring}, \code{noncomm\_ring [*]}, well-chosen
\code{simp} rules, and the same workers without information sharing. None of
those comparisons was executed here.

\subsection{Why this addition is useful}

Matrices, endomorphisms, projections, and operator expressions produce the same
obstruction: one knows equations about products, but the useful occurrence of a
relation is buried inside a larger product. For instance,
\[
 UA=1,\qquad BV=1
 \quad\Longrightarrow\quad
 U-V=U(B-A)V.
\]
An equality worker should find the two surrounding contexts instead of asking a
human to insert them. Likewise, $G\succeq0$ naturally proves $Q^*GQ\succeq0$,
not $GQ\succeq0$. And $\tr([A,B])=0$ is a scalar trace fact, not the assertion
$[A,B]=0$.

These are useful new outputs for a planner: an equality for local saturation,
a proved positive operator for an analytic subgoal, a scalar trace inequality,
or a concrete diagnostic against an over-general algebraic conjecture. They
need not replace any existing Forge worker.

\section{The object language and its exact interpretations}

\subsection{Words instead of exponent vectors}

Let $X=\{x_0,\ldots,x_{m-1}\}$ and let $X^*$ be its finite words, including
the empty word $\epsilon$. A polynomial in the free associative unital algebra
\[
 \F=\Q\langle x_0,\ldots,x_{m-1}\rangle
\]
is a finitely supported map $X^*\to\Q$. Multiplication concatenates words;
$xy$ and $yx$ remain distinct. The empty word represents $1$.

The implementation stores a word as a tuple of atom indices and a polynomial as
an exact sparse dictionary. Serialization uses sorted triples
\code{[word, numerator, denominator]}, with nonzero reduced coefficients and
positive denominators. The checker implements its own accumulation and product
routines; it does not import the producer's polynomial arithmetic or elimination
code. The representation and wire schema remain shared design choices, so this
is algorithmic separation, not protection against every common-mode mistake.

Rational coefficients are \emph{central}. A Lean source translation must target
an associative $\Q$-algebra, not an arbitrary ring. The implication $2x=0
\Rightarrow x=0$ admits the rational receipt $x=\tfrac12(2x)$ but is false in
characteristic two. Merely changing a pretty-printed scalar type would invalidate
the interface. Integer-coefficient receipts can be reconstructed more generally,
as several of the supplied Lean targets illustrate.

\subsection{Four meanings, not four names for the same object}

The problem and certificate schemas distinguish the following cases.

\begin{description}[style=nextline,leftmargin=0pt,labelwidth=0pt,labelsep=0pt]
\item[Algebra equality.]
For every associative unital $\Q$-algebra and every tuple satisfying the listed
relations $r_j=0$, the target $p$ is zero. The prototype schema intentionally
has no additional domain restrictions.
\item[Algebra countermodel.]
Concrete rational square matrices satisfy every listed relation as a whole
matrix, and a supplied vector is separated by $p$. This disproves the preceding
universal algebra equality. It does not refute a goal about a particular matrix
size, self-adjoint matrices, or a restricted algebra unless those extra
conditions are separately checked.
\item[Operator positivity.]
For real square matrices of arbitrary finite dimension, satisfying the declared
transpose involution on atoms, the relations, and the positive semidefinite
premises $g_i\succeq0$, the target satisfies $p\succeq0$.
\item[Trace positivity.]
Under the same matrix assumptions, $\tr(p)\ge0$. This is strictly weaker than
operator positivity and accepts a different residual language.
\end{description}

The checker requires the certificate kind to agree with the original problem,
except for the explicitly permitted equality/countermodel alternative. Unknown
fields such as an unsupported fixed matrix dimension or finite-field tag are
rejected rather than ignored. This does not solve the Lean source-translation
problem: a future reifier must first prove that the source goal has the
interpretation being requested.

\subsection{Involution and matrix semantics}

An involution on atom indices is a permutation $i\mapsto i^*$ with
$(i^*)^*=i$. Extend it by
\[
 (x_{i_1}\cdots x_{i_k})^*
   =x_{i_k^*}\cdots x_{i_1^*},\qquad c^*=c\quad(c\in\Q).
\]
Fixed atoms represent self-adjoint matrices. A two-cycle can represent a matrix
and its transpose. Both the reversal and the index map are indispensable. The
executed positivity corpus mostly uses fixed atoms; a separate unit test covers
paired atoms. The finite-ray producer currently assumes fixed atoms, whereas
the homogeneous Gram producer accepts either kind.

The operator and trace schemas require a formally self-adjoint target and
formally self-adjoint positive premises. Self-adjointness that follows only
\emph{modulo relations} requires an extra equality obligation or a rewritten
self-adjoint expression. This conservative entrance restriction caused a real
failed experiment in Section~\ref{sec:chsh}; it was preserved, not bypassed.

Complex matrices and abstract ordered star algebras are natural extensions of
the same mathematical proof rules. The delivered schema fixes real matrix
semantics; a larger interpretation deserves its own proved adapter rather than
an undocumented change of scope.
